% ============================================================
% section7_theory.tex - Раздел 7: Теоретическая модель проволочного поляризатора
% ============================================================
% !TeX program = xelatex
% Compartibility check: xelatex --shell-escape main.tex

\section{Теоретическая модель проволочного поляризатора}\label{sec:theory}

В предыдущих разделах мы научились загружать, визуализировать и предварительно обрабатывать экспериментальные данные ТГц-спектроскопии. Однако для глубокого анализа угловой зависимости пропускания нам необходима теоретическая модель. В этом разделе мы познакомимся с физической и математической основой работы металлических проволочных решеток в ТГц-диапазоне, опираясь на классические статьи A.~Blanco (1986) и K.~Kaltenecker (2019).

\subsection{Физический смысл проволочной решетки}

Идеальный поляризатор пропускает 100\% излучения с поляризацией, перпендикулярной проволокам, и полностью блокирует излучение с поляризацией, параллельной им. В реальности проволочные решетки ограничены геометрическими размерами~--- шагом решетки $p$ (периодом) и диаметром проволоки $d$. 

Когда длина волны падающего излучения $\lambda$ становится сопоставимой с шагом решетки $p$ (так называемая «дифракционная область», $p > \lambda$), эффективность поляризатора падает: появляется паразитное пропускание параллельной поляризации (обозначим его $t_\parallel$), а пропускание перпендикулярной ($t_\perp$) становится меньше единицы. Кроме того, возникают дифракционные эффекты и потери энергии на высшие моды, что требует более сложных математических моделей.

\subsection{Метод эквивалентных длинных линий (Transmission Line Analogy)}

Для описания взаимодействия электромагнитных волн с проволочной решеткой применяется мощный инженерный подход~--- аналогия с длинными линиями, предложенная Маркувицем (Marcuvitz). Согласно этой модели, электромагнитная волна, распространяющаяся в свободном пространстве, рассматривается как сигнал в линии передачи, а проволочная решетка~--- как препятствие с определенным комплексным импедансом (сопротивлением) $Z$.

При нормальном падении ТГц-волны выделяют два предельных случая:
\begin{itemize}
    \item \textbf{Емкостное препятствие (Capacitive obstacle):} Возникает, когда вектор электрического поля $\mathbf{E}$ \textbf{перпендикулярен} проволокам. Электроны в проволоках не могут свободно течь вдоль короткого направления (диаметра), и решетка ведет себя как распределенный конденсатор.
    \item \textbf{Индуктивное препятствие (Inductive obstacle):} Возникает, когда вектор электрического поля $\mathbf{E}$ \textbf{параллелен} проволокам. Электроны свободно разгоняются вдоль длинных металлических нитей, создавая сильный индукционный ток, что аналогично поведению распределенной катушки индуктивности.
\end{itemize}

Метод Бланко расширяет эту теорию на дифракционную область ($p > \lambda$), вводя комплексные импедансы с ненулевой действительной частью. Это позволяет учесть диссипацию энергии из-за связи с затухающими (эванесцентными) волнами высших порядков.

\subsection{Математический аппарат модели Бланко}

Приведем полную систему уравнений модели Бланко, связывающую геометрические параметры решетки ($p$, $d$) и длину волны излучения $\lambda$ с амплитудными коэффициентами пропускания $t_\perp$ и $t_\parallel$.

\subsubsection{Пропускание перпендикулярной составляющей (Capacitive case)}
Амплитудный коэффициент пропускания для перпендикулярной поляризации $t_\perp$ вычисляется как:
\begin{equation}
    t_\perp = \frac{2 Z_1 Z_2 / Z_0^2}{(1 + Z_1/Z_0)(Z_b/Z_0 + Z_2/Z_0)} \label{eq:t_perp}
\end{equation}
где $Z_0 \approx 377$~Ом~--- волновое сопротивление вакуума, а приведенные импедансы $Z_1/Z_0$ и $Z_2/Z_0$ выражаются через промежуточные компоненты $Z_a/Z_0$ и $Z_b/Z_0$:
\begin{align}
    \frac{Z_1}{Z_0} &= \frac{Z_a^2/Z_0^2 + Z_a Z_b / Z_0^2 + Z_a^2 Z_b / Z_0^3}{2 Z_a/Z_0 + Z_a^2/Z_0^2 + Z_b/Z_0 + Z_a Z_b / Z_0^2} \\
    \frac{Z_2}{Z_0} &= \frac{Z_a/Z_0}{1 + Z_a/Z_0}
\end{align}
Сами компоненты задаются чисто мнимыми величинами, зависящими от вспомогательных функций $f_a$ и $f_b$:
\begin{equation}
    \frac{Z_0}{Z_a} = i f_a\left(p, d, \lambda\right), \quad \frac{Z_0}{Z_b} = i f_b\left(p, d, \lambda\right)
\end{equation}
\begin{align}
    f_a(p, d, \lambda) &= \frac{p}{2\lambda} \left(\frac{\pi d}{p}\right)^2 \frac{1}{A_2} \label{eq:fa} \\
    f_b(p, d, \lambda) &= \frac{2\lambda}{p} \left(\frac{p}{\pi d}\right)^2 A_1 - \frac{p}{4\lambda} \left(\frac{\pi d}{p}\right)^2 \frac{1}{A_2} \label{eq:fb}
\end{align}

\subsubsection{Пропускание параллельной составляющей (Inductive case)}
Для параллельной поляризации амплитудный коэффициент пропускания $t_\parallel$ равен:
\begin{equation}
    t_\parallel = \frac{2 Z_3 Z_4 / Z_0^2}{(1 + Z_3/Z_0)(Z_d/Z_0 + Z_4/Z_0)(1 + Z_d/Z_0)} \label{eq:t_par}
\end{equation}
где приведенные импедансы $Z_3/Z_0$ и $Z_4/Z_0$ рассчитываются через реактивные компоненты $Z_c/Z_0$ и $Z_d/Z_0$:
\begin{align}
    \frac{Z_3}{Z_0} &= \frac{Z_d/Z_0 + 2 Z_c Z_d / Z_0^2 + Z_d^2 / Z_0^2 + Z_c/Z_0}{1 + Z_c/Z_0 + Z_d/Z_0} \\
    \frac{Z_4}{Z_0} &= \frac{Z_c/Z_0 + Z_c Z_d / Z_0^2}{1 + Z_c/Z_0 + Z_d/Z_0}
\end{align}
Импедансы определяются функциями $f_c$ и $f_d$:
\begin{equation}
    \frac{Z_c}{Z_0} = i f_c\left(p, d, \lambda\right), \quad \frac{Z_d}{Z_0} = -i f_d\left(p, d, \lambda\right)
\end{equation}
\begin{align}
    f_c(p, d, \lambda) &= \frac{p}{\lambda} \left[ \ln\left(\frac{p}{\pi d}\right) + \sum_{m=1}^{\infty} \left( \frac{1}{C_m} - \frac{1}{m} \right) \right] \label{eq:fc} \\
    f_d(p, d, \lambda) &= \frac{p}{\lambda} \left( \frac{\pi d}{p} \right)^2 \label{eq:fd}
\end{align}

\subsubsection{Геометрические коэффициенты связи \texorpdfstring{$A_1$, $A_2$ и $C_m$}{A1, A2 и Cm}}
Функции $A_1$ и $A_2$ учитывают ближнепольное взаимодействие между соседними нитями решетки:
\begin{align}
    A_1 &= 1 + \frac{1}{2} \left(\frac{\pi d}{\lambda}\right)^2 \left[ \ln\left(\frac{p}{\pi d}\right) + \frac{3}{4} \right] + \frac{1}{2} \left(\frac{\pi d}{\lambda}\right)^2 \sum_{m=1}^{\infty} \left( \frac{1}{C_m} - \frac{1}{m} \right) \label{eq:A1} \\
    A_2 &= 1 + \frac{1}{2} \left(\frac{\pi d}{\lambda}\right)^2 \left[ \frac{11}{4} - \ln\left(\frac{p}{\pi d}\right) \right] + \frac{1}{24}\left(\frac{\pi d}{p}\right)^2 - \left(\frac{\pi d}{p}\right)^2 \sum_{m=1}^{\infty} \left[ m - \frac{1}{2m}\left(\frac{p}{\lambda}\right)^2 - C_m \right] \label{eq:A2}
\end{align}
где коэффициент $C_m$ характеризует распространение $m$-й пространственной гармоники дифрагированного поля:
\begin{equation}
    C_m = \sqrt{m^2 - \left(\frac{p}{\lambda}\right)^2} \label{eq:Cm}
\end{equation}

При практических вычислениях бесконечные ряды в выражениях~(\ref{eq:fc}), (\ref{eq:A1}) и (\ref{eq:A2}) быстро сходятся, поэтому суммирование достаточно ограничить первыми 10--15 членами (в нашем коде принято $N=15$).

\note{Обрати внимание на коэффициент $C_m$. Когда $m < p/\lambda$, подкоренное выражение становится отрицательным, и $C_m$ приобретает комплексное значение. Физически это означает, что соответствующая пространственная гармоника перестает быть затухающей (эванесцентной) и начинает уносить энергию от решетки в виде распространяющегося дифракционного пучка.}

\subsection{Спектральные характеристики и физический анализ}

Чтобы наглядно понять поведение проволочного поляризатора в широком диапазоне частот, проанализируем спектральные зависимости амплитудных коэффициентов $|t_\perp|$ и $|t_\parallel|$, рассчитанные для паспортных параметров исследуемого в данной книге поляризатора ($p = 16.0$~мкм, $d = 11.0$~мкм).

\smartfigure{fig7.1.png}{Теоретические спектры амплитудного пропускания параллельной ($|t_\parallel|$) и перпендикулярной ($|t_\perp|$) составляющих проволочного поляризатора в диапазоне $0.1$--$100$~ТГц. Отчетливо виден глубокий резонансный провал (аномалия Вуда) на частоте $18.75$~ТГц.}{fig:7_1}

Данная спектральная зависимость (Рисунок~\ref{fig:7_1}) была сгенерирована с помощью вспомогательной функции \mycode{plot_polarizer_theory_characteristics()} из модуля визуализации \mycode{plots.py} (исходный код этой функции приведен в \textbf{\textit{Приложении~\ref{app:plots}}}). На графике отчетливо выделяются три характерные частотные области:

\begin{enumerate}
    \item \textbf{Рабочая длинноволновая область ($\nu \to 0$, $p \ll \lambda$):}
    На низких частотах (до $5$~ТГц) поляризатор работает идеально. Пропускание перпендикулярной составляющей $|t_\perp|$ практически равно $1.0$ (потери отсутствуют), а пропускание параллельной $|t_\parallel|$ стремится к нулю. Решетка эффективно разделяет поляризации.
    
    \item \textbf{Резонансная область (Аномалия Вуда, $\lambda \approx p$):}
    При приближении длины волны к периоду решетки возникает резкий электродинамический резонанс. На частоте $\nu_{\text{res}} = c/p \approx 18.75$~ТГц наблюдается глубокий провал пропускания $|t_\perp|$ и резкий всплеск паразитного пропускания $|t_\parallel|$. Этот эффект носит название \textbf{аномалии Вуда} (Wood's anomaly) и связан с тем, что первая дифрагированная волна начинает скользить вдоль поверхности решетки, эффективно перераспределяя энергию между модами.
    
    \item \textbf{Коротковолновая дифракционная область ($p > \lambda$):}
    На частотах выше $18.75$~ТГц решетка перестает работать как поляризатор и начинает вести себя как сложная дифракционная решетка. Коэффициенты $|t_\perp|$ и $|t_\parallel|$ испытывают хаотические осцилляции (резонансы высших порядков), а их значения сближаются, что делает невозможным использование решетки для поляризационных измерений.
\end{enumerate}

Таким образом, физически рабочий диапазон нашего поляризатора ограничен сверху частотой порядка $4$--$5$~ТГц, что полностью перекрывает спектр нашего терагерцового спектрометра (обычно работающего до $2.5$~ТГц).

\subsection{Допущения и границы применимости модели Бланко}

При использовании модели Бланко в научных и учебных расчетах необходимо четко понимать её ограничения:
\begin{enumerate}
    \item \textbf{Допущение тонких препятствий ($d \ll \lambda$ и $d \ll p$):}
    Математические формулы Бланко получены в предположении, что ток, наводимый в металлических проволоках, постоянен и однороден по их сечению. Это допущение отлично выполняется на низких частотах (в ТГц-области), где длина волны значительно больше диаметра нити ($\lambda \gg d$).
    
    \item \textbf{Случай «толстых» и близко расположенных решеток ($d/p \ge 0.5$):}
    В нашем реальном поляризаторе шаг равен $16.0$~мкм, а диаметр нити~--- $11.0$~мкм ($d/p \approx 0.69$). Из-за высокой плотности заполнения решетки и близости нитей друг к другу в окрестности резонанса ($\lambda \approx p$, т.е. около $18.75$~ТГц) электростатические приближения длинных линий начинают работать менее точно. Это может приводить к небольшим численным артефактам в модели (например, к флуктуациям суммарного баланса мощности вблизи точки сингулярности).
    
    \item \textbf{Численная регуляризация в точке сингулярности:}
    На частоте точного резонанса $p = \lambda$ коэффициент связи $C_1 = \sqrt{1^2 - (p/\lambda)^2}$ обращается в ноль, вызывая деление на ноль в бесконечных рядах Бланко. Для стабилизации расчетов в модуле \mycode{theoretical.py} применена малая мнимая регуляризация $C_m = \sqrt{m^2 - (p/\lambda)^2 + 10^{-9}i}$. Физически это эквивалентно введению бесконечно малых потерь (конечной проводимости проволок), что сглаживает математическую расходимость и делает алгоритм абсолютно устойчивым.
    
    \item \textbf{Рабочий терагерцовый диапазон:}
    В полезном диапазоне нашего ТГц-спектрометра (до $2.0$~ТГц) длина волны излучения составляет $\lambda \ge 150$~мкм. При этом геометрический параметр $d/\lambda \le 0.07$, что гарантирует высочайшую точность и физическую адекватность модели Бланко.
\end{enumerate}

\subsection{Матричный формализм Джонса и система двух поляризаторов}

В реальном эксперименте с THz-TDS (как описано в статье Kaltenecker, 2019) часто используется система из \textbf{двух} поляризаторов для контролируемого ослабления пучка. Первый поляризатор можно вращать на угол $\theta$, а второй жестко зафиксирован и пропускает только одну поляризацию (например, вдоль оси $x$).

\smartfigure{fig0.1.png}{Принципиальная схема эксперимента с двумя проволочными поляризаторами на пути ТГц излучения.}{fig:5_3}

Прохождение света через такую оптическую систему описывается матрицами Джонса. Поляризатор представляется диагональной матрицей $\mathbf{P}$, а его поворот~--- матрицей поворота $\mathbf{R}(\theta)$:
\begin{equation}
\mathbf{P} = \begin{pmatrix} t_\perp & 0 \\ 0 & t_\parallel \end{pmatrix}, \quad \mathbf{R}(\theta) = \begin{pmatrix} \cos\theta & -\sin\theta \\ \sin\theta & \cos\theta \end{pmatrix}
\end{equation}


Тогда электрическое поле на выходе из системы $\mathbf{E}_{out}$ связано со входным полем $\mathbf{E}_{in}$ (поляризованным по оси $x$) следующим образом:
\begin{equation}\label{eq:jones_system}
\mathbf{E}_{out} = \mathbf{P} \mathbf{R}(\theta) \mathbf{P} \mathbf{R}(-\theta) \mathbf{E}_{in}
\end{equation}
Здесь $\mathbf{R}(-\theta)$ переводит вектор поля в систему координат повернутого первого поляризатора, первая матрица $\mathbf{P}$ применяет пропускание, $\mathbf{R}(\theta)$ возвращает вектор в лабораторную систему координат, а вторая $\mathbf{P}$ описывает пропускание второго (неподвижного) поляризатора.

В результате аналитического перемножения этих матриц мы получаем выражение для выходного поля:
\begin{equation}
E_{out} = E_{in} \begin{pmatrix} t_\perp(t_\perp \cos^2\theta + t_\parallel \sin^2\theta) \\ t_\parallel(t_\perp - t_\parallel)\sin\theta\cos\theta \end{pmatrix}
\end{equation}

\subsection{Зачем нам эти формулы?}
В следующем разделе мы переведем эти громоздкие формулы в элегантный Python-код. Модель Бланко позволит нам:
\begin{enumerate}
    \item Вычислять $T_\perp(\lambda)$ и $T_\parallel(\lambda)$ для любых параметров решетки $p$ и $d$.
    \item Описывать деполяризацию пучка и возникновение эллиптичности.
    \item Строить теоретические кривые пропускания мощности $P(\theta, \lambda)$ и напрямую сравнивать их с нашими экспериментальными данными из датасета.
\end{enumerate}

Таким образом, мы сможем не просто наблюдать графики, а проверять, насколько хорошо реальный терагерцовый поляризатор описывается классической электродинамической моделью эквивалентных длинных линий.
